\documentclass[11pt,letterpaper]{article}
\usepackage[utf8]{inputenc}
\usepackage[T1]{fontenc}
\usepackage[spanish]{babel}
\usepackage[margin=2cm]{geometry}
\usepackage{enumitem}
\usepackage{booktabs}
\usepackage{amsmath}
\usepackage{hyperref}
\usepackage{listings}
\usepackage{xcolor}
\definecolor{codebg}{rgb}{0.95,0.95,0.95}
\lstset{basicstyle=\ttfamily\small,backgroundcolor=\color{codebg},frame=single,breaklines=true}

\title{Manual de Presentaci\'on\\Gu\'ia para Defensa de Tesis:\\Dise\~no y Construcci\'on de un Exoesqueleto\\Optimizando los Mecanismos para Agarre de Pinza Fina}
\author{Ing. Adriana Elorza Ramos}
\date{\today}

\begin{document}
\maketitle
\tableofcontents
\newpage

\section{Introducci\'on - Qu\'e decir ante los sinodales}

\subsection{Contexto cl\'inico}
La pinza fina es la funci\'on motora que permite al pulgar e \'indice sujetar objetos peque\~nos con precisi\'on. Su deterioro, causado frecuentemente por accidentes cerebrovasculares (ACV), genera p\'erdida de independencia en actividades como escribir, comer o vestirse.

\subsection{Justificaci\'on del exoesqueleto}
Los exoesqueletos permiten ejercicios de rehabilitaci\'on repetitivos y precisos, reduciendo la carga del terapeuta. La terapia rob\'otica ha demostrado ser tan eficaz como la convencional (Fukui 2025, Ueki 2012).

\subsection{Alcance del proyecto}
Exoesqueleto de eslabones r\'igidos para los tres dedos principales del agarre de pinza fina (\'indice, medio, pulgar), con movimiento en las tres articulaciones: MCP (metacarpofalangica), PIP (interfalangica proximal) y DIP (interfalangica distal).

\subsection{Por qu\'e eslabones r\'igidos}
\begin{itemize}
\item Control bidireccional (asiste flexi\'on Y extensi\'on)
\item Predecibles cinem\'aticamente (modelo anal\'itico exacto)
\item Permiten optimizaci\'on dimensional directa
\item Los cables tienen fricci\'on variable y requieren tensores adicionales
\end{itemize}

\section{Objetivo General}

Dise\~nar y construir un exoesqueleto de eslabones r\'igidos para realizar agarres de pinza fina, empleando t\'ecnicas de an\'alisis cinem\'atico y de optimizaci\'on para adaptar los par\'ametros del mecanismo a distintas medidas antropom\'etricas de la poblaci\'on mexicana.

\subsection{Justificaci\'on ante sinodales}
El movimiento del dedo humano es no lineal, resultado de la interacci\'on coordinada de m\'ultiples articulaciones. Un dise\~no por tanteo no puede reproducir esta complejidad. La optimizaci\'on permite explorar sistem\'aticamente el espacio de dise\~no (18 variables) para minimizar el error entre la trayectoria del mecanismo y la trayectoria anat\'omica real.

\section{Trabajo Anterior}

\subsection{Modelo CAD}
Dise\~no 3D completo contemplando tres mecanismos (uno por dedo), sistemas de transmisi\'on (engranes, tornillo sin f\'in) y motores.

\subsection{Selecci\'on de componentes}
Se evaluaron 6 tipos de motores mediante una matriz de decisi\'on ponderada con 7 criterios: torque (20\%), velocidad (20\%), tama\~no (20\%), costo (10\%), \'angulo de giro (10\%), peso (10\%), voltaje (10\%). El micro-reductor Pololu N20 obtuvo 4.0/5.0 por su torque de 6.7 kgf$\cdot$cm y tama\~no compacto (13.6$\times$42.5 mm).

\subsection{Construcci\'on}
Piezas fabricadas mediante impresi\'on 3D (PLA/PETG). Verificaci\'on de ensamble y movimiento de cada mecanismo de forma independiente.

\section{Base de Datos}

\subsection{Dataset de Pizzolato et al. (2017)}
Base de datos calibrada publicada en Scientific Data con cinem\'atica y EMG de 40 sujetos realizando diferentes tipos de agarres. Captura con CyberGlove (guante sensorizado) y marcadores \'opticos.

\subsection{Extracci\'on de datos}
Se seleccion\'o el agarre de pinza fina (125$^\circ$ de flexi\'on total). Se extrajeron los \'angulos articulares del dedo \'indice (MCP, PIP, DIP) y se calcul\'o la cinem\'atica directa:
\begin{align}
px_{IFP} &= F_p \cos(\theta_{MCP}) \\
py_{IFP} &= F_p \sin(\theta_{MCP}) \\
px_{IFD} &= px_{IFP} + F_m \cos(\theta_{MCP} + \theta_{PIP}) \\
py_{IFD} &= py_{IFP} + F_m \sin(\theta_{MCP} + \theta_{PIP})
\end{align}

Medidas antropom\'etricas (poblaci\'on mexicana adulta): $F_p = 49$ mm, $F_m = 26$ mm, $F_d = 24$ mm.

\section{Evoluci\'on del Algoritmo de Optimizaci\'on}

\subsection{Fase 1: AG + Hiperheur\'istica}

\textbf{Qu\'e es una hiperheur\'istica:} Una heur\'istica que selecciona o genera heur\'isticas. Decide din\'amicamente qu\'e operador gen\'etico usar: mutaci\'on suave o agresiva, cruce aritm\'etico o uniforme.

\textbf{Configuraci\'on:}
\begin{itemize}
\item 12 variables de dise\~no
\item Modelo cinem\'atico: 5 barras + 4 barras $\rightarrow$ posici\'on IFP
\item Poblaci\'on: 60, Generaciones: 120, Mutaci\'on: 25\%, Cruce: 80\%
\item Selecci\'on: torneo binario, Elitismo global
\end{itemize}

\textbf{Resultado:} Error RMS = 0.006 m (6 mm) --- insuficiente.

\textbf{Problemas:}
\begin{enumerate}
\item Alta sensibilidad: ecuaciones de cierre frecuentemente sin soluci\'on
\item Eslabones no realistas: longitudes $>$ 10 cm
\item Costo computacional excesivo
\item Hiperheur\'istica insuficiente para equilibrar trayectoria y restricciones
\end{enumerate}

\subsection{Fase 2: Transici\'on a Evoluci\'on Diferencial (DE)}

\textbf{Por qu\'e DE supera a AG (Price, Storn \& Lampinen, 2005):}

\begin{tabular}{lll}
\toprule
Aspecto & AG & DE \\
\midrule
Representaci\'on & Binaria o real & Vectores reales directos \\
Mutaci\'on & Perturbaci\'on fija & Auto-adaptativa (vectores diferencia) \\
Convergencia (18 var) & Lenta & Probada eficiente \\
Paisaje discontinuo & Sensible & Robusto \\
\bottomrule
\end{tabular}

\vspace{0.3cm}
\textbf{C\'omo funciona DE:}
$$\vec{v}_i = \vec{x}_{r1} + F \cdot (\vec{x}_{r2} - \vec{x}_{r3})$$

El tama\~no del paso se adapta autom\'aticamente a la escala del problema.

\subsection{Fase 3: Optuna (Optimizaci\'on Bayesiana)}

Framework de optimizaci\'on bayesiana (Akiba et al., KDD 2019) que usa TPE (Tree-structured Parzen Estimator).

\textbf{Estrategia de 2 etapas:}
\begin{enumerate}
\item Optuna ejecuta 8 trials cortos (20 iteraciones c/u) buscando mejor configuraci\'on de DE (estrategia, popsize, mutaci\'on, recombinaci\'on)
\item Con la mejor configuraci\'on, DE ejecuta 1000 iteraciones + polish con L-BFGS-B
\end{enumerate}

\subsection{Fase 4: Distancia de Chamfer}

\textbf{Por qu\'e no RMS:} El RMS punto-a-punto asume correspondencia ordenada. Si hay desfase temporal, el RMS se infla aunque las formas sean similares.

\textbf{Distancia de Chamfer (Borgefors, 1988):}
$$d_{CH}(S,T) = \frac{1}{|S|}\sum_{s \in S} \min_{t \in T} \|s - t\| + \frac{1}{|T|}\sum_{t \in T} \min_{s \in S} \|t - s\|$$

Ventajas: bidireccional, robusta ante diferentes densidades, tolerante a desalineamientos.

\subsection{Fase 5: Correcci\'on de ``ganchos'' (exo\_14 a exo\_17)}
\begin{enumerate}
\item Inversi\'on temporal: datos MOCAP ven\'ian de flexi\'on m\'axima $\rightarrow$ extensi\'on; se invirtieron
\item Recorte DIP: \'angulos negativos (hiperextensi\'on por ruido) recortados a 0$^\circ$
\item Suavizado Savitzky-Golay: filtro polinomial (ventana 15, orden 2)
\item Penalizaci\'on de monotonicidad: si $\theta_{fm}$ retrocede $\rightarrow$ soluci\'on descartada
\end{enumerate}

\section{Explicaci\'on Detallada del C\'odigo (exo\_17.py)}

\subsection{Estructura general}
507 l\'ineas divididas en 9 secciones: par\'ametros, carga MOCAP, funciones de evaluaci\'on, soluciones matem\'aticas, modelo cinem\'atico, funci\'on objetivo, bounds, Optuna, ejecuci\'on principal.

\subsection{RESTRICCIONES (11 tipos)}

\begin{enumerate}
\item \textbf{gear\_ratio $>$ 0:} Evita divisi\'on por cero en $\theta_1 = \theta_2/\text{gear\_ratio}$
\item \textbf{min(p[:12]) $>$ 0.005:} Fabricabilidad m\'inima de 5 mm (l\'imite de impresi\'on 3D)
\item \textbf{hsp $>$ 0 y dsp $>$ 0:} Los soportes deben tener dimensi\'on positiva f\'isica
\item \textbf{Cierre 1er 5 barras:} Discriminante $g^2 - 4dh \geq 0$ y denominador $\neq 0$. Si no se cumple, las barras no alcanzan a unirse.
\item \textbf{Cierre 1er 4 barras:} Discriminante $B_1^2 - 4A_1C_1 \geq 0$. Misma raz\'on.
\item \textbf{Cierre 2do 5 barras:} Misma restricci\'on para el segundo lazo del mecanismo.
\item \textbf{Cierre 2do 4 barras:} Misma restricci\'on para el segundo lazo.
\item \textbf{Anti-gancho:} $\Delta\theta_{fm} \geq -0.5^\circ$ entre pasos consecutivos. Un dedo real no retrocede durante el cierre.
\item \textbf{Espacio de trabajo:} $|px_{tip}| < 0.3$ m, $|py_{tip}| < 0.3$ m, sin NaN. Valores fuera indican geometr\'ia degenerada.
\item \textbf{Bounds:} 18 variables dentro de l\'imites (manejado internamente por scipy DE).
\item \textbf{Penalizaci\'on de monotonicidad:} $W_{MONO} = 5.0$ aplicada a IFD y TIP en la funci\'on fitness.
\end{enumerate}

\textbf{Por qu\'e existen:}
\begin{itemize}
\item Restricciones de cierre (4--7): Si no se cumplen, el mecanismo \textbf{no puede ensamblarse} f\'isicamente. Las barras no alcanzan a unirse.
\item Monotonicidad (8, 11): Un dedo real no retrocede durante el cierre. Una trayectoria con ``ganchos'' es f\'isicamente imposible.
\item Fabricabilidad (2, 3): Eslabones menores a 5 mm no se pueden fabricar con precisi\'on.
\item Bounds (10): Dimensiones mayores a 120 mm exceden el tama\~no de la mano humana.
\end{itemize}

\subsection{PESOS en la funci\'on objetivo}

\begin{lstlisting}
W_IFP  = 0.25    # Peso articulacion IFP
W_IFD  = 0.375   # Peso articulacion IFD
W_TIP  = 0.375   # Peso punta del dedo
W_MONO = 5.0     # Penalizacion monotonicidad
\end{lstlisting}

\textbf{Por qu\'e IFP tiene MENOS peso (0.25):}
\begin{itemize}
\item La articulaci\'on IFP tiene trayectoria relativamente simple (arco circular)
\item Es la m\'as cercana al actuador, m\'as f\'acil de replicar
\item En versiones intermedias ya converg\'ia bien; el problema era IFD y TIP
\end{itemize}

\textbf{Por qu\'e IFD y TIP tienen M\'AS peso (0.375):}
\begin{itemize}
\item Son articulaciones distales donde se acumulan errores de toda la cadena
\item Presentaban los ``ganchos'' (inversiones de trayectoria)
\item Para pinza fina, la posici\'on de la PUNTA es lo m\'as cr\'itico (hace contacto con el objeto)
\end{itemize}

\textbf{Por qu\'e $W_{MONO} = 5.0$:}
\begin{itemize}
\item Peso alto porque ganchos son F\'ISICAMENTE IMPOSIBLES
\item Determinado experimentalmente: $< 3.0$ aparec\'ian ganchos; $> 8.0$ convergencia muy lenta
\item Valor 5.0 es el balance \'optimo entre suavidad y velocidad de convergencia
\end{itemize}

\subsection{Modelo cinem\'atico paso a paso}

Para cada \'angulo de entrada $\theta_2$ (de 0$^\circ$ a 85$^\circ$, 120 puntos):
\begin{enumerate}
\item C\'alculo de $\theta_1 = \theta_2 / \text{gear\_ratio} + \text{offset}$
\item \textbf{1er 5 barras:} Resuelve punto P mediante ecuaci\'on cuadr\'atica (intersecci\'on de circunferencias)
\item \textbf{1er 4 barras:} Calcula $\theta_4$ mediante m\'etodo de semi-tangente $\rightarrow$ posici\'on IFP
\item \textbf{Soportes:} Calcula posiciones de soporte de falange proximal ($c_1$, $\theta_{ps1}$, $rs_2$, $\theta_{aux,s2}$)
\item \textbf{2do 5 barras:} En sistema de referencia rotado $\rightarrow$ punto P2
\item \textbf{2do 4 barras:} Obtiene $\theta_{fm}$, $\theta_{fd}$ (\'angulos de falanges media y distal)
\item \textbf{Posiciones cartesianas:} $IFD = F_m \cos(\theta_{fm}) + IFP$, $TIP = F_d \cos(\theta_{fd}) + IFD$
\item \textbf{Verificaci\'on:} anti-gancho y espacio de trabajo
\end{enumerate}

\subsection{Funci\'on fitness}
\begin{enumerate}
\item Ejecuta \texttt{run\_kinematics} (si retorna None $\rightarrow$ fitness = 1000.0)
\item Concatena puntos MOCAP y simulados (IFP + IFD + TIP)
\item Procrustes: encuentra R, t \'optimos mediante SVD
\item Aplica transformaci\'on r\'igida a puntos simulados
\item Calcula Chamfer distance por articulaci\'on
\item Calcula penalizaci\'on de monotonicidad en IFD y TIP
\item Retorna: $0.25 \cdot err_{IFP} + 0.375 \cdot err_{IFD} + 0.375 \cdot err_{TIP} + 5.0 \cdot (mono_{IFD} + mono_{TIP})$
\end{enumerate}

\section{Resultados Finales}

\subsection{Par\'ametros optimizados}

\begin{tabular}{clcc}
\toprule
\# & Par\'ametro & Valor & Bounds \\
\midrule
1 & Bancada1 & 71.18 mm & [15, 120] mm \\
2 & Bancada2 & 53.86 mm & [15, 120] mm \\
3 & Link1 & 17.68 mm & [10, 120] mm \\
4 & Link2 & 81.94 mm & [10, 120] mm \\
5 & Link3 & 21.43 mm & [10, 120] mm \\
6 & Link4 & 67.01 mm & [10, 120] mm \\
7 & Link5 & 50.29 mm & [10, 120] mm \\
8 & Link6 & 23.68 mm & [10, 120] mm \\
9 & Link7 & 34.60 mm & [10, 120] mm \\
10 & Link8 & 80.38 mm & [10, 120] mm \\
11 & Link9 & 48.27 mm & [10, 120] mm \\
12 & Link10 & 64.50 mm & [10, 120] mm \\
13 & hsp & 57.28 mm & [-40, 80] mm \\
14 & dsp & 24.49 mm & [5, 80] mm \\
15 & $\theta_{aux,FM}$ & 0.8931 rad & [0, $\pi$] \\
16 & $\theta_{aux,FD}$ & 0.3582 rad & [0, $\pi$] \\
17 & Rel. engranaje & 1.7281 & [1.0, 8.0] \\
18 & $\theta_{offset}$ & $-$2.0149 rad & [$-\pi$, $\pi$] \\
\bottomrule
\end{tabular}

\subsection{Comparaci\'on de resultados}
\begin{tabular}{lcc}
\toprule
M\'etrica & AG + Hiperheur\'istica & DE + Optuna \\
\midrule
Error global & 6 mm & $<$ 1 mm \\
Eslabones & No realistas ($>$10 cm) & Dentro de bounds \\
Ganchos & Presentes & Eliminados \\
\bottomrule
\end{tabular}

\section{Preguntas Frecuentes de Sinodales}

\begin{enumerate}
\item \textbf{?`Por qu\'e 18 variables?} El mecanismo compuesto (2 lazos 5B + 2 lazos 4B) necesita todas para describir completamente la geometr\'ia y el acoplamiento.
\item \textbf{?`C\'omo garantizas \'optimo global?} No se garantiza en metaheur\'istica. Sin embargo, m\'ultiples configuraciones de Optuna + polish L-BFGS-B dan buena cobertura. El error sub-milim\'etrico est\'a cerca del l\'imite del ruido del sensor.
\item \textbf{?`Para qu\'e la relaci\'on de engranaje?} Permite que un motor con rango limitado cubra el rango completo de flexi\'on del dedo (0--85$^\circ$ de entrada $\rightarrow$ 0--49$^\circ$ de MCP).
\item \textbf{?`Qu\'e pasa si cambias medidas antropom\'etricas?} Se re-ejecuta el algoritmo con nuevos $F_p$, $F_m$, $F_d$. Esa es la ventaja principal del enfoque de optimizaci\'on.
\item \textbf{?`Por qu\'e Chamfer y no Fr\'echet?} Chamfer tiene costo $O(nm)$ vs $O(n^2m^2)$ de Fr\'echet. Con Procrustes + monotonicidad, Chamfer da resultados comparables.
\end{enumerate}

\section{Trabajo Futuro}
\begin{enumerate}
\item Armar modelo CAD con las 18 dimensiones optimizadas
\item Validar estudio de movimiento en software CAD vs trayectorias MOCAP
\item Reconstruir prototipo f\'isico con nuevas medidas (impresi\'on 3D)
\item Repetir optimizaci\'on para dedos medio y pulgar
\item Pruebas funcionales del exoesqueleto completo
\end{enumerate}

\end{document}
